% bpc.tex — Branch-Price-and-Cut 求解 CVRP 子问题
% 符号体系与 lbbd.tex / reformulationModel.tex / speed.tex 保持一致

\documentclass{article}
\usepackage{amsmath,amssymb,amsfonts,amsthm}
\usepackage{enumitem}
\usepackage{geometry}
\usepackage{mathrsfs}
\geometry{a4paper, margin=1in}
\usepackage{ctex}
\usepackage{hyperref}
\usepackage{float}
\usepackage[linesnumbered,ruled,vlined]{algorithm2e}

\newtheorem{remark}{注}
\newtheorem{definition}{定义}
\newtheorem{proposition}{命题}

\begin{document}

\title{基于自适应大邻域搜索的启发式求解方法}
\maketitle  % 必须添加这一行，LaTeX 才会把声明的标题渲染出来
\label{chap:alns}

上一章建立的 Original Model 能够通过 CPLEX 等求解器获得精确最优解，
但随着收集点数量 $n$、时期数 $l$ 和车辆数 $K$ 的增长，
混合整数规划（MIP）的求解时间迅速膨胀，难以在可接受的时间内求解中大规模实例。
因此，本章设计一种基于自适应大邻域搜索（Adaptive Large Neighbourhood Search, ALNS）
的元启发式算法来高效求解 MVPRP。

%% ============================================================
\section{算法框架选择的理论依据}
\label{sec:alns-motivation}
%% ============================================================

\subsection{为何选择 Original Model 作为解空间的基础}

第~\ref{chap:model}~章同时介绍了 Original Model 与基于最短路重构的
Reformulation Model。后者通过引入 $\lambda_{ivt}$ 变量
将收集点库存约束编码进最短路网络，使线性规划（LP）松弛更紧，
适用于列生成与对偶割的精确求解框架。
然而在启发式环境中，算法直接操作可行解而非求解 LP 松弛，
Reformulation 在松弛质量上的优势无法发挥作用。

Original Model 的决策变量——访问标记 $z_{it}$、取货量 $q_{it}$、
路径弧 $x_{ijt}$、生产量 $p_t$——均具有直接的物理含义。
邻域操作（如插入或移除一次访问、交换路径弧段、调整生产量）
的定义与可行性检查因此更为直观。
相比之下，$\lambda_{ivt}$ 的最短路流平衡约束
在每次邻域移动后都需要维护链式一致性，
增加实现复杂度却无法在启发式框架中带来解质量的提升。
因此，本章选择 Original Model 的变量空间作为 ALNS 的操作基础。

\subsection{为何选择 ALNS 框架}

MVPRP 耦合了三个子结构：批量生产决策（lot-sizing）、
多层库存管理（inventory management）、以及带容量约束的多车路径规划（CVRP）。
ALNS 的"破坏—修复"（destroy–repair）框架天然适合此类多维决策问题，
原因有三：

\begin{enumerate}
  \item \textbf{子结构专用算子。}
    可以针对不同子结构设计专门的破坏与修复算子：
    "移除某收集点在某时期的访问"作用于库存层面；
    "重建某时期的路径"作用于路径层面；
    "扰动一段时间窗内的生产计划"作用于生产层面。
    这些算子可独立设计、灵活组合。
  \item \textbf{自适应权重机制。}
    ALNS 通过自适应权重调整自动学习当前实例上最有效的算子组合，
    比固定邻域结构（如禁忌搜索）更具鲁棒性与灵活性。
  \item \textbf{文献基础成熟。}
    在库存路径问题（IRP）与生产路径问题（PRP）领域，
    ALNS 已被广泛采用并积累了大量成熟的算子设计经验
    \cite{Ropke2006,Pisinger2007,Coelho2014,Adulyasak2015}。
\end{enumerate}

%% ============================================================
\section{解的表示与评价}
\label{sec:alns-representation}
%% ============================================================

\subsection{解的编码}

一个完整解 $\mathcal{S}$ 由以下组件构成：

\begin{itemize}
  \item \textbf{访问计划} $z_{it} \in \{0,1\}$：
    收集点 $i$ 在时期 $t$ 是否被访问，$i=1,\dots,n$，$t=1,\dots,l$。
  \item \textbf{取货量} $q_{it} \geq 0$：
    若 $z_{it}=1$，则 $q_{it}$ 等于访问时收集点 $i$ 的全部库存
    （即清空策略）；否则 $q_{it}=0$。
  \item \textbf{路径方案} $\mathcal{R}_t$：
    时期 $t$ 的一组车辆路径，每条路径为一个收集点序列，
    满足载重约束 $Q$ 与车辆数约束 $K$。
  \item \textbf{生产计划} $(y_t, p_t, P^0_t, I^0_t)$：
    分别为时期 $t$ 的生产启动标志、生产量、成品库存和原料库存。
\end{itemize}

关键的层次依赖关系为：$z$ 与 $\mathcal{R}$ 确定后，
$q$ 由清空策略直接计算；
给定 $q$，工厂原料库存 $I^0$ 由库存平衡方程确定；
生产计划 $(y_t, p_t, P^0_t)$ 则由独立的子问题优化得出。

\subsection{目标函数分解}

总目标函数可分解为三个部分：
\begin{equation}
  Z(\mathcal{S}) = Z_{\text{route}} + Z_{\text{inv}} + Z_{\text{prod}}
  \label{eq:alns-objective}
\end{equation}
其中：
\begin{align}
  Z_{\text{route}}
    &= \sum_{t=1}^{l} \sum_{(i,j) \in \mathcal{R}_t} c_{ij}
    \label{eq:alns-route-cost} \\
  Z_{\text{inv}}
    &= \sum_{t=1}^{l} \sum_{i=1}^{n} h_i \cdot I_{it}
    \label{eq:alns-holding-cost} \\
  Z_{\text{prod}}
    &= \sum_{t=1}^{l} \bigl( u \cdot p_t + f \cdot y_t
       + h_0 \cdot I^0_t + h_p \cdot P^0_t \bigr)
    \label{eq:alns-prod-cost}
\end{align}
其中 $I_{it}$ 为收集点 $i$ 在时期 $t$ 末的库存，
$c_{ij}$ 为从节点 $i$ 到节点 $j$ 的运输成本。

\subsection{可行性检查}

每次破坏-修复操作后，需验证以下硬约束：

\begin{enumerate}
  \item \textbf{收集点库存容量：}
    对所有 $i,t$，取货前库存
    $I_{i,t-1} + s_{it} \leq L_i$。
  \item \textbf{工厂原料与成品库存容量：}
    $I^0_t \leq L_0$ 且 $P^0_t \leq L_p$。
  \item \textbf{车辆数约束：}
    每期使用的路径数 $|\mathcal{R}_t| \leq K$。
  \item \textbf{载重约束：}
    每条路径的总载重 $\sum_{i \in r} q_{it} \leq Q$。
\end{enumerate}

%% ============================================================
\section{增量评价机制}
\label{sec:alns-incremental}
%% ============================================================

ALNS 的核心效率瓶颈在于每次候选移动的评价成本。
若每次移动都执行完整的解拷贝与全量重新评价，计算代价将十分高昂。
因此，本文设计了一套增量评价（incremental evaluation）机制，
仅计算移动引起的目标函数变化量 $\Delta Z$，
而无需重建整个解。

\subsection{收集点库存的增量计算}

对于收集点 $i$，其库存轨迹完全由访问序列 $\{t : z_{it}=1\}$ 决定。
当考虑翻转某个 $z_{it}$ 时，只需沿时间轴重算该收集点的库存序列，
时间复杂度为 $O(l)$。具体做法为：临时翻转 $z_{it}$ 的值，
沿 $t=1,\dots,l$ 顺序扫描该收集点的库存演化，
计算新的持有成本与溢出状态，然后恢复原值。
该过程不产生任何对象分配，全部在原始解上原地完成。

\subsection{路径成本的增量计算}

对于单次插入或移除操作，路径成本变化量可由邻接节点的距离差直接计算。
设收集点 $c$ 位于路径中位置 $\text{pos}$，
其前驱为 $\text{prev}$，后继为 $\text{next}$（depot 以节点 $0$ 和 $n{+}1$ 表示），则：

\begin{itemize}
  \item 移除 $c$ 的路径成本变化：
    $\Delta_{\text{remove}} = c_{\text{prev},\text{next}}
     - c_{\text{prev},c} - c_{c,\text{next}}$
  \item 在 $\text{prev}$ 与 $\text{next}$ 之间插入 $c$：
    $\Delta_{\text{insert}} = c_{\text{prev},c} + c_{c,\text{next}}
     - c_{\text{prev},\text{next}}$
\end{itemize}

对于同路径内的 relocate 或 swap 操作，
由于前驱-后继关系发生耦合变化，
本文采用完整的路径成本重算（复杂度 $O(|\text{route}|)$）
以保证正确性。

\subsection{综合移动评价}

一个访问移除的综合评价 $\Delta_{\text{total}}$ 由路径成本变化量
$\Delta_{\text{route}}$ 与库存持有成本变化量 $\Delta_{\text{inv}}$ 之和构成：
\begin{equation}
  \Delta_{\text{total}} = \Delta_{\text{route}} + \Delta_{\text{inv}}
  \label{eq:incremental-total}
\end{equation}
同时，增量评价还返回布尔标志指示移动是否导致收集点库存溢出，
供后续的可行性过滤使用。

%% ============================================================
\section{初始解的构造}
\label{sec:alns-initial}
%% ============================================================

良好的初始解能显著加速 ALNS 的收敛。
本文采用"截止期驱动访问规划 $+$ 节约算法构路"的两阶段策略。

\subsection{截止期驱动的访问计划}

利用 $\mu(i,t)$ 参数的语义——
收集点 $i$ 在时期 $t$ 被最后一次访问后，
能维持到的最晚时期——
逐期构建访问计划：

\begin{enumerate}
  \item 对每个时期 $t$，首先插入所有满足
    $\mu(i, \text{prev}_i) = t$ 的强制访问
    （即若不在 $t$ 时期访问收集点 $i$，
    则后续必然发生库存溢出）。
  \item 检查当前累计原料收集量是否满足生产下界需求。
    若不足，贪心选择能提供最大取货量且不违反容量约束的收集点加入。
  \item 每次插入前，通过增量评价验证插入不会导致
    该收集点在当期或更早时期发生库存溢出。
\end{enumerate}

\subsection{路径构造}

对每个时期 $t$，采用 Clarke-Wright 节约算法
\cite{Clarke1964} 构造初始路径：

\begin{enumerate}
  \item 初始化：每个被访问的收集点单独构成一条路径。
  \item 计算所有收集点对 $(i,j)$ 的节约值
    $s_{ij} = c_{0i} + c_{j,n{+}1} - c_{ij}$。
  \item 按节约值降序排列，依次尝试合并路径尾部为 $i$
    与头部为 $j$ 的两条路径，
    前提是合并后载重不超过 $Q$。
  \item 若合并后路径数仍超过 $K$，
    退化为贪心装箱策略：
    按取货量降序排列收集点，
    依次将每个收集点插入现有路径中成本增量最小的位置。
\end{enumerate}

\subsection{生产计划优化}

给定访问计划与路径确定的取货量 $q_{it}$ 后，
生产计划 $(y_t, p_t, P^0_t, I^0_t)$ 构成一个确定性的
单产品批量生产问题（lot-sizing problem），
其变量规模仅为 $O(l)$，
可通过求解如下混合整数规划快速得到最优解：
\begin{align}
  \min \quad & \sum_{t=1}^{l} \bigl( u \cdot p_t + f \cdot y_t
    + h_0 \cdot I^0_t + h_p \cdot P^0_t \bigr) \\
  \text{s.t.} \quad
    & P^0_{t-1} + p_t = d_t + P^0_t
      & \forall\, t \label{eq:prod-finished-balance} \\
    & p_t \leq C \cdot y_t
      & \forall\, t \label{eq:prod-capacity} \\
    & I^0_{t-1} + \textstyle\sum_{i} q_{it} = k \cdot p_t + I^0_t
      & \forall\, t \label{eq:prod-raw-balance} \\
    & I^0_t \leq L_0,\quad P^0_t \leq L_p
      & \forall\, t \\
    & y_t \in \{0,1\},\quad p_t, P^0_t, I^0_t \geq 0
      & \forall\, t
\end{align}
由于该子问题规模极小，CPLEX 通常在毫秒级内返回最优解。
这一设计将生产决策从 ALNS 的邻域搜索中解耦，
使路径与库存层面的算子无需关心生产可行性的维护。

%% ============================================================
\section{破坏算子}
\label{sec:alns-destroy}
%% ============================================================

本文设计了六种破坏算子，分别从随机性、成本敏感性、
空间相关性、时间结构和生产耦合等不同维度破坏当前解。
每次破坏操作移除的访问数量 $d$ 由以下公式自适应确定：
\begin{equation}
  d = \text{rand}\bigl[1, \;
      \min\bigl(|\mathcal{V}|, \;
      \max(d_{\min}, \lceil \rho \cdot n \cdot l \rceil)\bigr)\bigr]
  \label{eq:destroy-count}
\end{equation}
其中 $|\mathcal{V}|$ 为当前解中的总访问数，
$d_{\min}$ 为最小破坏量，$\rho$ 为破坏比例参数。

\subsection{D1: 随机访问移除（Random Visit Removal）}

从当前解的所有访问 $(i,t)$ 中均匀随机选取 $d$ 个，
将其从解中移除。
该算子提供最大的随机性，确保搜索空间的多样性。

\subsection{D2: 最差成本移除（Worst Cost Removal）}

对每个访问 $(i,t)$，利用增量评价计算其边际成本影响：
\begin{equation}
  \text{score}(i,t) = \Delta_{\text{route}}(i,t) + \Delta_{\text{inv}}(i,t)
  \label{eq:worst-cost-score}
\end{equation}
其中 $\Delta_{\text{route}}$ 为移除该访问带来的路径成本减少量，
$\Delta_{\text{inv}}$ 为移除后库存持有成本的变化量（通常为正，
因为库存会在该期继续累积）。
按 $\text{score}$ 升序排列（即边际贡献最差的访问排在前面），
选取前 $d$ 个移除。

对于移除后会导致库存溢出的访问，施加惩罚项以降低其被移除的优先级：
\begin{equation}
  \text{score}'(i,t) = \begin{cases}
    \text{score}(i,t) + \alpha \cdot \bar{\sigma}
      & \text{若移除导致溢出} \\
    \text{score}(i,t) & \text{否则}
  \end{cases}
  \label{eq:worst-cost-penalty}
\end{equation}
其中 $\bar{\sigma}$ 为所有非溢出访问的平均边际成本绝对值，
$\alpha > 0$ 为惩罚因子。

\subsection{D3: 相关访问移除（Related Visit Removal）}

随机选择一个种子访问 $(i_0, t_0)$，
计算所有其他访问与种子的相关度：
\begin{equation}
  \text{relatedness}(i,t;\; i_0, t_0)
    = (1-\omega) \cdot \frac{c_{i_0 i}}{c_{\max}}
      + \omega \cdot \frac{|t_0 - t|}{l - 1}
  \label{eq:relatedness}
\end{equation}
其中 $c_{\max} = \max_{i,j} c_{ij}$ 为最大距离，
$\omega \in [0,1]$ 为时间权重参数，用于控制空间相似性与时间相近性之间的平衡。
按相关度升序选取前 $d$ 个访问移除。

直觉在于：地理位置接近且时间相邻的访问群体被同时移除后，
留出一个连贯的"空白区域"，便于修复阶段系统性地重新规划。

\subsection{D4: 时期破坏（Period Destruction）}

随机选取 $1$ 至 $2$ 个时期，完全清除这些时期的所有访问与路径。
该算子提供最大粒度的破坏，适用于某些时期的规划质量整体不佳的情况。

\subsection{D5: 路径移除（Route Removal）}

随机选取某个时期中的一条完整路径，
移除该路径上的所有访问。
该算子的破坏范围介于 D1 和 D4 之间，
专注于路径结构层面的重组。

\subsection{D6: 生产引导扰动（Production-Guided Perturbation）}

随机选取 $1$ 至 $2$ 个当前有生产活动的时期，
将其生产标志 $y_t$ 和生产量 $p_t$ 置零，
同时在相邻时期以一定概率移除部分访问。
该算子迫使修复阶段重新寻找原料供应与生产分配的平衡，
是唯一直接作用于生产子结构的破坏算子。

%% ============================================================
\section{修复算子}
\label{sec:alns-repair}
%% ============================================================

\subsection{修复任务的识别}

修复阶段首先识别当前部分解中的两类不可行性来源，
并将其抽象为待处理的"修复任务"：

\begin{enumerate}
  \item \textbf{收集点溢出任务：}
    若收集点 $i$ 在某时期 $t^*$ 发生库存溢出
    （即 $I_{i,t^*-1} + s_{it^*} > L_i$），
    则生成一个任务 $(i, t^*)$，要求在 $t^*$ 之前（含 $t^*$）
    增加对 $i$ 的访问。
  \item \textbf{原料不足任务：}
    若累计收集的原料不足以支撑生产下界要求，
    则生成一个任务，要求在对应时期之前增加原料收集量。
\end{enumerate}

所有任务按截止期升序排列，确保最紧迫的需求优先满足。

\subsection{R1: 贪心插入修复（Greedy Insertion Repair）}

对任务列表按截止期顺序处理，
对每个任务枚举所有候选插入位置
（收集点 $\times$ 时期 $\times$ 路径位置），
选择总增量成本 $\Delta_{\text{total}}$ 最小的位置执行插入。
候选过滤条件包括：

\begin{itemize}
  \item 插入后路径载重不超过 $Q$；
  \item 时期总载重不超过 $K \cdot Q$；
  \item 插入后该收集点不在截止期内发生新的溢出。
\end{itemize}

\subsection{R2: 遗憾值插入修复（Regret Insertion Repair）}

遗憾值插入 \cite{Ropke2006} 旨在避免贪心插入的近视性。
对每个待处理任务，计算其"遗憾值"——
最优插入位置与次优插入位置的成本差：
\begin{equation}
  \text{regret}(\tau) = \Delta^{(2)}_\tau - \Delta^{(1)}_\tau
  \label{eq:regret}
\end{equation}
其中 $\Delta^{(k)}_\tau$ 表示任务 $\tau$ 的第 $k$ 好插入成本。
优先处理遗憾值最大的任务——直觉在于，
遗憾值大的任务若延迟处理，其好位置可能被占据，
导致后续只能选择代价高得多的位置。

\subsection{R3: 逐期重建修复（Period-by-Period Reconstruction）}

对受破坏影响的时期（包括被 D4 完全清除的时期
和被其他算子部分移除访问的时期），
按时间顺序逐期重建：

\begin{enumerate}
  \item 首先插入截止期约束强制要求的访问（$\mu(i, \text{prev}_i) = t$ 的收集点）；
  \item 若该期原料收集量不足，贪心添加取货量最大的收集点；
  \item 使用 Clarke-Wright 节约算法重建该期路径。
\end{enumerate}

重建完成后，仍可能存在跨期的可行性缺口，
此时调用贪心插入（R1）进行兜底修复。

\subsection{R4: 生产计划再优化}

每次修复完成后，\textbf{始终}调用第~\ref{sec:alns-initial}~节
所述的批量生产子问题求解器，
重新优化生产计划。
这保证了在路径与库存决策变更后，
生产计划总是关于当前取货方案的最优解。

%% ============================================================
\section{局部搜索}
\label{sec:alns-local-search}
%% ============================================================

每次修复与生产再优化完成后，
对当前解执行一轮局部搜索以进一步提升解质量。
局部搜索分为两个阶段依次执行。

\subsection{跨时期局部搜索}

\textbf{跨期 relocate} 操作尝试将收集点 $i$ 的一次访问
从时期 $t_1$ 移至时期 $t_2$。
该操作同时涉及两个时期的路径变更与全局库存的重新计算。

由于跨期移动的精确评价需要调用生产再优化，计算代价较高，
本文采用"估算-筛选-精算"的三阶段策略：

\begin{enumerate}
  \item \textbf{估算阶段：}
    对每个候选移动，基于路径成本变化与库存持有成本变化
    估算目标函数改善量，同时通过快速的生产可行性检验
    （基于原料累计量的上下界包络线）过滤显然不可行的移动。
  \item \textbf{筛选阶段：}
    按估算改善量排序，仅保留前 $M$ 个候选。
  \item \textbf{精算阶段：}
    对筛选通过的候选执行完整的解构造（含生产再优化），
    接受首个确认改善的移动。
\end{enumerate}

其中，快速生产可行性检验通过维护一个成品库存的
可达区间 $[\underline{P}^0_t, \overline{P}^0_t]$ 实现，
判断条件为：若存在某时期使得 $\underline{P}^0_t > \overline{P}^0_t$，
则该取货方案必然不可行。

\subsection{期内局部搜索}

对每个时期内的路径执行以下三种经典 VRP 局部搜索算子，
直到无法继续改善为止：

\begin{enumerate}
  \item \textbf{2-opt：}
    对路径中的弧段 $(i, i{+}1, \dots, k, k{+}1)$
    尝试反转为 $(k, k{-}1, \dots, i{+}1, i)$，
    若路径成本减少则执行。

  \item \textbf{路径间 relocate：}
    将一个收集点从路径 $r_1$ 移至路径 $r_2$ 的最优位置，
    须满足移动后 $r_2$ 的载重不超过 $Q$。

  \item \textbf{路径间 swap：}
    交换两条路径中各一个收集点的位置，
    须满足交换后两条路径的载重均不超过 $Q$。
    对于同一路径内的 swap，无需进行容量检查
    （因为总载重不变），仅验证路径成本改善。
\end{enumerate}

%% ============================================================
\section{ALNS 主循环}
\label{sec:alns-main-loop}
%% ============================================================

\subsection{模拟退火接受准则}

采用模拟退火（Simulated Annealing）作为接受准则。
设当前解的目标函数值为 $Z_{\text{curr}}$，
候选解为 $Z_{\text{cand}}$，则候选解以如下概率被接受：
\begin{equation}
  P(\text{accept}) = \begin{cases}
    1 & \text{若 } Z_{\text{cand}} \leq Z_{\text{curr}} \\
    \exp\!\Bigl(-\dfrac{Z_{\text{cand}} - Z_{\text{curr}}}{T}\Bigr)
      & \text{否则}
  \end{cases}
  \label{eq:sa-acceptance}
\end{equation}

温度 $T$ 按运行时间进度做指数衰减：
\begin{equation}
  T = T_0 \cdot \gamma^{\,\text{progress}}
  \label{eq:temperature}
\end{equation}
其中 $\text{progress} = \frac{\text{已用时间}}{\text{总时限}} \in [0,1]$，
$\gamma$ 为终止温度比率（$\gamma = T_{\text{final}} / T_0$）。
初始温度设定为：
\begin{equation}
  T_0 = \frac{0.05 \cdot Z_{\text{init}}}{\ln 2}
  \label{eq:initial-temperature}
\end{equation}
使得初始状态下目标值恶化 $5\%$ 的解有约 $50\%$ 的接受概率。

使用基于时间而非迭代次数的温度衰减，
使得在不同规模实例上（单次迭代耗时差异可能很大）
温度曲线保持一致的搜索行为。

\subsection{自适应权重更新}

每个破坏算子 $d$ 和修复算子 $r$ 分别维护权重 $w_d$ 和 $w_r$。
算子通过轮盘赌选择：算子 $d$ 被选中的概率为
$w_d / \sum_{d'} w_{d'}$。

每隔一个 segment（$\eta$ 次迭代），
根据算子在该 segment 内的表现更新权重。
算子在每次使用时获得如下得分：
\begin{equation}
  \text{score} = \begin{cases}
    \sigma_1 & \text{若产生新的全局最优解} \\
    \sigma_2 & \text{若优于当前解（但非全局最优）} \\
    \sigma_3 & \text{若被接受但劣于当前解} \\
    0        & \text{若不被接受}
  \end{cases}
  \label{eq:alns-score}
\end{equation}

权重更新公式为：
\begin{equation}
  w_d \leftarrow (1 - r) \cdot w_d + r \cdot \frac{\text{Score}_d}{\text{Uses}_d}
  \label{eq:weight-update}
\end{equation}
其中 $r \in (0,1)$ 为反应系数。
为防止权重趋于零导致某些算子永远不被选中，
设置下界 $w_{\min} = 10^{-3}$。

\subsection{重启机制}

当搜索陷入停滞时，执行重启（restart）操作。
停滞的判定采用双重条件：全局最优解连续
$\theta_{\text{best}}$ 次迭代未改善，
\textbf{且}当前解连续 $\theta_{\text{curr}}$ 次迭代未改善。
双重条件的目的是避免在全局最优未更新但当前解仍在有效探索
（通过接受劣解实现多样化）时过早重启。

重启操作包括：
\begin{enumerate}
  \item 将当前解回退为历史全局最优解；
  \item 在随后的 $\eta$ 次迭代中，
    使用更大的破坏比例 $\rho_{\text{restart}} > \rho$
    和更高的温度乘数 $\beta > 1$，
    以增强搜索多样性。
\end{enumerate}

\subsection{算法伪代码}

完整的 ALNS 主循环如算法~\ref{alg:alns-main} 所示。

\begin{algorithm}[htbp]
\caption{ALNS 求解 MVPRP 的主循环}
\label{alg:alns-main}
\begin{algorithmic}[1]
\REQUIRE 问题实例 $\mathcal{I}$，参数集合 $\Theta$
\ENSURE 最优可行解 $\mathcal{S}^*$
\STATE $\mathcal{S} \leftarrow \textsc{ConstructInitialSolution}(\mathcal{I})$
\STATE $\mathcal{S}^* \leftarrow \mathcal{S}$
\STATE $T_0 \leftarrow 0.05 \cdot Z(\mathcal{S}) / \ln 2$
\STATE 初始化所有算子权重 $w_d = w_r = 1$
\STATE $\text{iter} \leftarrow 0$
\WHILE{未超过时间限制}
  \STATE $\text{iter} \leftarrow \text{iter} + 1$
  \IF{满足重启条件}
    \STATE $\mathcal{S} \leftarrow \mathcal{S}^*$；
    进入 reheat 模式（增大破坏比例，升高温度乘数）
  \ENDIF
  \STATE 以轮盘赌选择破坏算子 $d$ 与修复算子 $r$
  \STATE $\mathcal{S}' \leftarrow \textsc{Copy}(\mathcal{S})$
  \STATE $\textsc{Destroy}_d(\mathcal{S}')$
  \STATE $\textsc{Repair}_r(\mathcal{S}')$
  \STATE $\textsc{ProductionReoptimize}(\mathcal{S}')$
  \STATE $\textsc{LocalSearch}(\mathcal{S}')$
  \IF{$\mathcal{S}'$ 可行}
    \IF{以模拟退火准则接受 $\mathcal{S}'$}
      \STATE $\mathcal{S} \leftarrow \mathcal{S}'$
    \ENDIF
    \IF{$Z(\mathcal{S}') < Z(\mathcal{S}^*)$}
      \STATE $\mathcal{S}^* \leftarrow \mathcal{S}'$
    \ENDIF
  \ENDIF
  \STATE 更新算子得分
  \IF{$\text{iter} \bmod \eta = 0$}
    \STATE 更新自适应权重
  \ENDIF
\ENDWHILE
\RETURN $\mathcal{S}^*$
\end{algorithmic}
\end{algorithm}

%% ============================================================
\section{参数设置}
\label{sec:alns-parameters}
%% ============================================================

表~\ref{tab:alns-params} 汇总了 ALNS 算法的主要参数及其取值。

\begin{table}[htbp]
\centering
\caption{ALNS 算法参数设置}
\label{tab:alns-params}
\begin{tabular}{llll}
\toprule
\textbf{参数} & \textbf{符号} & \textbf{取值} & \textbf{说明} \\
\midrule
segment 大小      & $\eta$         & 100   & 权重更新周期 \\
反应系数          & $r$            & 0.2   & 权重更新的学习率 \\
最优得分          & $\sigma_1$     & 10    & 产生全局最优解的奖励 \\
改善得分          & $\sigma_2$     & 5     & 优于当前解的奖励 \\
接受得分          & $\sigma_3$     & 1     & 接受劣解的奖励 \\
终止温度比率      & $\gamma$       & $10^{-3}$ & $T_{\text{final}}/T_0$ \\
最小破坏量        & $d_{\min}$     & 6     & 每次破坏的最少访问数 \\
破坏比例          & $\rho$         & 0.25  & 常规破坏比例 \\
重启破坏比例      & $\rho_{\text{restart}}$ & 0.35 & 重启后的增大破坏比例 \\
重启温度乘数      & $\beta$        & 3.0   & 重启后的温度放大系数 \\
溢出惩罚因子      & $\alpha$       & 1.5   & D2 中的溢出惩罚 \\
时间权重          & $\omega$       & 0.6   & D3 中的时间相关度权重 \\
重启停滞阈值（全局）& $\theta_{\text{best}}$ & $10\eta$ & 全局最优停滞触发 \\
重启停滞阈值（当前）& $\theta_{\text{curr}}$ & $3\eta$  & 当前解停滞触发 \\
期内搜索预算      & ---            & 40    & 每轮期内局部搜索的最大移动数 \\
跨期搜索预算      & ---            & 20    & 每轮跨期局部搜索的最大移动数 \\
\bottomrule
\end{tabular}
\end{table}

%% ============================================================
\section{本章小结}
\label{sec:alns-summary}
%% ============================================================

本章针对 MVPRP 设计了一种 ALNS 元启发式算法。
算法的核心设计思路可归纳为以下几点：

\begin{enumerate}
  \item \textbf{子问题解耦：}
    通过将生产决策封装为独立的批量优化子问题，
    ALNS 的邻域搜索仅需关注访问计划与路径，
    降低了搜索空间的维度。

  \item \textbf{多样化与集中化的平衡：}
    六种破坏算子从不同维度提供搜索多样性，
    自适应权重机制动态调整各算子的使用频率，
    局部搜索则在每次迭代中对当前解进行集中化改善。

  \item \textbf{增量评价驱动效率：}
    通过对库存成本和路径成本的增量计算，
    避免了频繁的全量解评价，使单次迭代的计算代价保持在较低水平。

  \item \textbf{双重停滞检测与自适应重启：}
    重启机制结合增大的破坏范围与升高的温度，
    有效避免搜索过早收敛到局部最优。
\end{enumerate}

下一章将在标准测试算例上对比 ALNS 与精确求解器的结果，
验证算法的有效性和效率。

\end{document}